\section{Conclusiones}
\label{iv:conclusiones}

El recorrido desarrollado enlaza la generación aritmética, la incidencia
excepcional y la realización dimensional mediante operaciones
especificadas. Su antecedente es el estado APP--TRIT--TPK con
orientación, hojas, cocientes, frontera y memoria. Las coordenadas
numéricas y los objetos excepcionales resultan de representaciones
de ese estado. Conservar el antecedente permite explicar conjuntamente
qué se evalúa y qué estructura hace posible la evaluación.

El registro \(K\) proporciona un resultado central de esta exposición.
Su construcción dodecafásica precede tanto a la lectura racional de
sus bloques como a la selección de incidencia y a la dirección
\(u_K\). En la representación real fijada, la positividad de
\(\|u_K\|^2\) determina una recta y el proyector de rango \(10\).
En el corredor excepcional, los soportes y la bandera seleccionada
conducen a la construcción del retículo y a la realización de
\(V^\natural\). La misma procedencia no hace idénticos estos
codominios: proporciona los mapas concretos que permiten compararlos
conservando sus acciones y la información que transmiten.

La representación de fase añade una consecuencia operatoria precisa
a la conservación de la unidad. La prolongación unital de todas las
fases conserva la traza normalizada y produce el cierre UHF nonádico.
Su representación tracial es un factor hiperfinito de tipo \(II_1\).
Los rangos finitos normalizados admiten en él todas las dimensiones
traciales del intervalo unidad. La continuación marcada tiene otra
inclusión y otro cierre. Este contraste identifica la operación que
conserva la normalización global y muestra por qué el dato de
refinamiento forma parte del resultado.

La dualidad de radio tiene aquí una formulación escalar y otra
operatoria, unidas por la misma involución. La igualdad de formas
cuadráticas se prolonga a una conjugación unitaria de Hamiltonianos
autoadjuntos con sus dominios. El cambio de sección residual requiere
transportar la relación \(L_9\); sobre ese cociente, la evaluación
mínima y la dualidad covariante están bien definidas. Las identidades
se mantienen para radios recíprocos y admiten la normalización de
escala de cuerda expuesta en el texto.

Las reducciones \(12\to11\to10\) y \(26\to24\) conservan,
respectivamente, la dirección determinada por \(K\) y la
presentación isotrópica del retículo. El morfismo conjunto de pantallas
coordina ambas con el sector compacto. Su prueba identifica los
cocientes, los núcleos y las transformaciones entrelazadas. La
perforación ternaria añade una realización combinatoria con su
distancia y su partición perfecta, cuyo rango queda explicado por
la operación sobre el código.

La acción sobre trayectorias se construye conservando el dato
precedente en cada composición. Los cambios de dirección y de tipo
aritmético forman entonces un cociclo exactamente aditivo. Ese
resultado fija una propiedad interna del transporte genealógico.
La superficie de mundo, la acción de Polyakov y los campos
undecadimensionales ocupan la etapa de realización: su compatibilidad
se expresa mediante los mapas y las condiciones operatorias
declarados. El teorema de transporte de la supercarga conserva
\(Q^2=H\) sobre la imagen correspondiente y precisa los dominios
necesarios para esa igualdad.

El contenido demostrado reúne, por tanto, la incidencia que alcanza
Moonshine, la dirección dodecafásica, la realización hilbertiana,
las pantallas, la dualidad compacta y la composición de acción.
La identificación física completa conserva los requisitos de sus
campos, tensiones, osciladores, amplitudes y límites declarados.
La distinción permite expresar con exactitud la contribución de cada
construcción: las pruebas algebraicas proporcionan una estructura
determinada y la realización física establece qué dinámica conserva
esa estructura. Este es el criterio de continuidad que enlaza el
presente artículo con los desarrollos posteriores de estadística,
radiación y estructura de las especies.
